\clearpage
\fancyhead[L]{\small CHRONOMETRISCHE LÄNGENBESTIMMUNGEN.}
\hypertarget{native-460-chronometric}{}
\expandafter\def\csname HistoricalNoteText462\endcsname{[beobachtet von Herrn Doctor \textsc{Tiarks} bei der von der Englischen Admiralität im Jahre 1824
veranstalteten Expedition zur chronometrischen Bestimmung der Längsunterschiede zwischen Altona,
Bremen, Helgoland und Greenwich. \textsc{Schering}.]}
\begin{center}{\large CHRONOMETRISCHE LÄNGENBESTIMMUNGEN.}\end{center}

\begin{center}
\rule{0.55\linewidth}{0.4pt}\\[2pt]
Astronomische Nachrichten. Band V. Nr. 110. Seite 227\,..\,234. 1826 November.\\[2pt]
\rule{0.55\linewidth}{0.4pt}
\end{center}

Es seien $\theta, \theta', \theta''\ldots$ die Zeiten (zusammen an der Zahl $n$), wo der Chronometer
vor den Zeiten der Oerter, deren Längen $x, x', x''$ etc. sind, um die Unterschiede
$a, a', a''$ etc. voraus war. Die Angaben $\theta, \theta', \theta''$ etc. setze ich schon auf Einen Ort
reducirt voraus. Ist also der tägliche Gang des Chronometers $=u$, so würde man, wenn
der Chronometer vollkommen wäre, die $n-1$ Gleichungen haben.

\[
a-\theta u-x = a'-\theta' u-x' = a''-\theta'' u-x'' = a'''-\theta''' u-x''' = \text{etc.}
\]

\begingroup\rightskip=0pt plus 1em
Damit diese Gleichungen zureichen, um die unbekannten Grössen $u, x, x', x''$ etc. zu
bestimmen, wird theils eine der Grössen $x, x', x''$ etc. als gegeben angesehen, theils
vorausgesetzt, dass wenigstens an Einem Orte zweimal beobachtet ist, also zwei der
Grössen $x, x', x''$ etc. identisch sind. Falls nun nicht mehr als zwei identisch sind, wird
die Aufgabe ganz bestimmt sein. Im entgegengesetzten Fall ist sie überbestimmt; und man
wird dann die unbekannten Grössen so bestimmen müssen, dass den $n-1$ Gleichungen\par\endgroup

\[
\begin{aligned}
0 &= a-a'+(\theta'-\theta)u-x+x'\\
0 &= a'-a''+(\theta''-\theta')u-x'+x''\\
0 &= a''-a'''+(\theta'''-\theta'')u-x''+x''' \quad \text{etc.}
\end{aligned}
\]
\noindent so genau wie möglich Genüge geleistet werde, da die immer stattfindenden
Unvollkommenheiten aller Chronometer nicht verstatten werden, Allen genau Genüge zu
leisten. Offenbar aber darf diesen Gleichungen nicht gleiches Gewicht beigelegt werden;
denn in der That drücken die Grössen

\[
\begin{aligned}
&a-a'+(\theta'-\theta)u-x+x'\\
&a'-a''+(\theta''-\theta')u-x'+x''
\end{aligned}
\]
\noindent bloss die Aggregate aller Abweichungen vom mittlern Gange aus, die der Chronometer in
den Zwischenzeiten $\theta'-\theta$, $\theta''-\theta'$ u.\,s.\,w. gehabt hat, und wenn von einem guten
Chronometer die Rede ist, dem man wirklich einen mittlern, keinen allmählig in einerlei
Sinn zunehmenden Aenderungen unterworfenen Gang beilegen kann, so wird der
mittlere zu befürchtende Werth eines solchen Aggregats der Quadratwurzel der
Zwischenzeit proportional gesetzt werden müssen.

\Needspace{3\baselineskip}
\SourceLeaf{461}
Demzufolge wird man also den obigen Gleichungen, indem man sie den Vorschriften der
Methode der kleinsten Quadrate gemäss behandelt, ungleiche Gewichte, die den
Zwischenzeiten $\theta'-\theta$, $\theta''-\theta'$, $\theta'''-\theta''$ etc. umgekehrt proportional sind,
beilegen müssen.

Die Auflösung hat dann keine Schwierigkeit, und man erhält sowohl die plausibelsten
Werthe von $u, x, x', x''$ etc. als ihre relative Zuverlässigkeit. Hiebei mache ich noch ein
Paar Bemerkungen.

1) Wenn die erste und letzte Beobachtung an Einerlei Orte gemacht sind, so ist der
plausibelste Werth von $u$ genau derselbe, der bloss aus der Vergleichung der beiden
äussersten Beobachtungen folgt. Die Rechnung wird dann ausserordentlich einfach, da es
nach einem leicht zu beweisenden Lehrsatz erlaubt ist, diesen plausibelsten Werth von
$u$ sogleich in den Gleichungen zu substituiren, oder, was dasselbe ist, die sämmtlichen
beobachteten Chronometerzeiten auf die eines fingirten zu reduciren, dessen Voreilung
$=0$ wäre.

2) Hat man den Gleichungen schlechtweg die Gewichte $\dfrac{1}{\theta'-\theta}$, $\dfrac{1}{\theta''-\theta'}$
beigelegt, so liegt den Gewichten, welche man für die Endresultate der
Längenbestimmungen findet, als Einheit die Genauigkeit zum Grunde, die man mit
\emph{diesem} Chronometer zu erwarten hätte, wenn man, bei bekanntem Gange, einen
Längenunterschied nach einem Zeitintervall von einem Tage bestimmte (insofern die
Zeiten $\theta, \theta', \theta''$ in Tagen ausgedrückt sind). Allein damit man die Resultate
\emph{verschiedener} Chronometer von ungleicher Güte vergleichen kann, muss noch ein
Factor hinzukommen, der von der Güte jedes einzelnen Chronometers abhängig ist.
Diesen zu finden, setze man die Werthe der Grössen

\[
\begin{aligned}
&a-a'+(\theta'-\theta)u-x+x'\\
&a'-a''+(\theta''-\theta')u-x'+x''\\
&a''-a'''+(\theta'''-\theta'')u-x''+x''' \quad \text{etc.}
\end{aligned}
\]
\noindent indem man für $u, x, x', x''$ etc. die gefundenen plausibelsten Werthe substituirt,
$=\lambda, \lambda', \lambda''$ etc. und

\[
\frac{\lambda\lambda}{\theta'-\theta} + \frac{\lambda'\lambda'}{\theta''-\theta'} + \frac{\lambda''\lambda''}{\theta'''-\theta''} + \text{etc.} = S
\]

Es sei ferner $\nu$ die Anzahl der sämmtlichen unbekannt gewesenen Grössen, und
$m=\BareRoot{\dfrac{S'}{n-\nu-1}}$, dann ist jener specifische Factor für jeden einzelnen Chronometer
der Grösse $\dfrac{1}{mm}$ oder $\dfrac{n-\nu-1}{S}$ proportional. Man kann $m$ als die mittlere zu
befürchtende Abweichung vom mittlern Gange nach Einem Tage Zwischenzeit ansehen.

3) Die obigen Vorschriften gelten für einen Chronometer, der keine erhebliche
progressive Abänderung seines Ganges zeigt. Wo das Gegentheil eintritt, kann man,
insofern die Reihe der Beobachtungen nicht übermässig lang ist, sich damit begnügen,
eine der Zeit proportionirte Abänderung des täglichen Ganges anzunehmen, so dass noch
eine unbekannte Grösse mehr einzuführen ist und die Gleichungen diese Gestalt haben:

\[
\begin{aligned}
0 &= a-a'+(\theta'-\theta)u+(\theta'\theta'-\theta\theta)v-x+x'\\
0 &= a'-a''+(\theta''-\theta')u+(\theta''\theta''-\theta'\theta')v-x'+x''\\
0 &= a''-a'''+(\theta'''-\theta'')u+(\theta'''\theta'''-\theta''\theta'')v-x''+x'''\\
&\text{u.\,s.\,w.}
\end{aligned}
\]

4) Dieses noch weiter zu treiben und also noch eine unbekannte Grösse mehr und
Glieder der Form $(\theta'^3-\theta^3)w$ einzuführen, möchte kaum rathsam sein. Chronometer, die
starke entschiedene Abänderungen des mittlern Ganges zeigen, die aber selbst wieder
unregelmässig sind, würde ich, neben andern, lieber ganz ausschliessen, da ihre
Resultate theils viel weniger genau werden, theils die Ge\SourceLeaf{462}nauigkeit sich viel schwerer in Zahlen zur Vergleichung angeben lässt. Ich halte mich daher hier bei
der viel verwickeltern Theorie solcher Fälle nicht auf, da in dem vorliegenden Fall das obige zureichend
sein wird.

4) Was die Auflösung der Gleichungen nach der Methode der kleinsten Quadrate betrifft, so ist
vielleicht nicht überflüssig in Erinnerung zu bringen, dass man in den meisten Fällen wohl thut, die
unbekannten Grössen aus einem bekannten (möglich genäherten) und einem unbekannten (also sehr
kleinen) Theile zusammenzusetzen. Dieser Rath ist zwar theils sonst schon wiederholt gegeben, theils
ist der Vortheil dieser Manier von selbst einleuchtend, allein es schien gut, ihn wieder in Erinnerung
zu bringen, da ich sehe, dass er häufig vergessen wird, wodurch die numerischen Rechnungen unnöthigerweise erschwert und Fehler leichter möglich werden.

Von den 36 Chronometern\HistoricalNote{462} habe ich folgende 5 berechnet.

\begin{center}
\begingroup\fontsize{8}{10}\selectfont\setlength{\tabcolsep}{2pt}
\begin{tabular}{ll|r|r|r|r|r|r}
 & & & Nr. 1 & Nr. 4 & \textsc{Breguet} 3056 & \textsc{Kessels} 1252 & \textsc{Barraud} 904 \\
\hline
Greenwich & Juni 30 & $3^{\rm u}22^{\rm m}$ & $-8^{\rm m}17^{\rm s}14$ & $+1^{\rm m}2^{\rm s}37$ &  &  &  \\
 & Juli 25 & $2\;15$ & $10\;44.39$ & $1\;32.15$ & $+30^{\rm m}59^{\rm s}75$ & $+50^{\rm m}29^{\rm s}31$ & $+48^{\rm m}29^{\rm s}20$ \\
 & 28 & $3\;13$ & $11\;0.69$ & $1\;36.96$ & $30\;50.07$ & $50\;39.69$ & $48\;40.24$ \\
 & Aug. 2 & $1\;15$ & $11\;28.48$ & $1\;44.44$ & $30\;31.78$ & $50\;52.14$ & $48\;58.87$ \\
 & 17 & $10\;28$ & $12\;59.40$ & $2\;6.24$ & $29\;35.69$ & $51\;38.66$ & $49\;57.83$ \\
 & 25 & $7\;27$ & $13\;47.98$ & $2\;15.84$ & $29\;10.48$ & $52\;2.45$ & $50\;27.15$ \\
 & Sept. 10 & $7\;40$ & $15\;24.47$ & $2\;40.36$ &  &  &  \\
Helgoland & Juli 3 & $3\;40$ & $-40\;8.00$ & $-30\;26.84$ &  &  &  \\
 & 22 & $12\;40$ & $42\;2.02$ & $30\;3.89$ & $-0\;20.34$ & $+18\;48.39$ & $+16\;47.39$ \\
 & Aug. 5 & $1\;48$ & $43\;18.11$ & $29\;43.35$ & $1\;10.24$ & $19\;26.77$ & $17\;37.51$ \\
 & 11 & $13\;9$ & $43\;35.77$ & $29\;33.43$ & $1\;32.75$ & $19\;47.22$ & $18\;1.30$ \\
 & 30 & $19\;30$ & $45\;53.08$ & $29\;7.96$ & $2\;40.67$ & $20\;47.68$ & $19\;17.03$ \\
 & Sept. 6 & $3\;6$ & $46\;31.56$ & $28\;58.94$ & $3\;4.55$ & $21\;6.56$ & $19\;43.80$ \\
 & 7 & $8\;42$ & $46\;38.72$ & $28\;56.71$ &  &  &  \\
Altona & Aug. 6 & $5\;55$ & $-51\;38.95$ & $-37\;55.76$ & $-9\;28.50$ & $+11\;16.25$ & $+9\;28.48$ \\
 & 9 & $12\;35$ & $51\;57.35$ & $37\;50.03$ & $9\;38.81$ & $11\;27.76$ & $9\;40.30$ \\
 & 31 & $9\;57$ & $54\;10.33$ & $37\;21.30$ & $10\;56.68$ & $12\;35.96$ & $11\;5.92$ \\
 & Sept. 4 & $22\;12$ & $54\;39.16$ & $37\;15.21$ & $11\;15.36$ & $12\;48.10$ & $11\;24.49$ \\
Bremen & Aug. 13 & $0\;2$ & $-47\;50.65$ & $-33\;16.49$ & $-5\;23.37$ & $+16\;5.83$ & $+14\;21.86$ \\
\end{tabular}
\endgroup
\end{center}

Ich setze die Rechnung für \emph{Breguet} 3056 zur Probe her. Die Länge von Helgoland sei $=0$, die
von Greenwich $=-x$, die von Altona $=+y$; Bremen schliesse ich hier aus, da es ohnehin, weil nur
einmal daselbst beobachtet ist, keine Controlle darbietet.

Die Zeiten rechne ich von der ersten Vergleichung der englischen Chronometer an (Greenwich
Jun. 30. $3^{\rm u}22^{\rm m}$).

Ich finde so, indem ich einen fingirten Chronometer vom Gange $=0$ substituire, dessen Stand.

\begin{center}
\begin{tabular}{r|l||r|l}
$\theta$ & & $\theta$ & \\
\hline
22.4 & $+60^{\rm s}20$ & 42.4 & $+59^{\rm s}88$ \\
25.0 & $+1949.60-x$ & 48.3 & $+1949.60-x$ \\
28.0 & $+1950.87-x$ & 56.2 & $+1952.74-x$ \\
32.9 & $+1950.29-x$ & 61.6 & $+61.32$ \\
35.9 & $+59.08$ & 62.2 & $-432.53+y$ \\
37.1 & $-434.98+y$ & 66.8 & $-434.98+y$ \\
40.4 & $-433.49+y$ & 68.0 & $+60.19$ \\
\end{tabular}
\end{center}

\SourceLeaf{463}
Die obigen Gleichungen fallen nun hier so aus, dass $x$ und $y$ gar nicht gemischt sind, wodurch
die weitere Rechnung noch bequemer wird. Wir haben nemlich für $x$ vier Bestimmungen.

\[
\begin{array}{ll}
+1889^{\rm s}40 & \text{Gewicht } \dfrac{1}{2.6}=0.38 \\[7pt]
+1891.21 & \dfrac{1}{3.0}=0.33 \\[7pt]
+1889.78 & \dfrac{1}{5.9}=0.17 \\[7pt]
+1891.42 & \dfrac{1}{3.4}=0.19 \\[7pt]
\hline
\text{Also } x=+1890^{\rm s}36 \quad \text{Gewicht} & =1.07
\end{array}
\]

Ebenso findet man

\[
y=+494^{\rm s}12 \quad \text{Gewicht } =3.83
\]
\noindent substituirt man diese Werthe, so ist der Stand des fingirten Chronometers gegen Helgolander Zeit

\begin{center}
\begingroup\small\renewcommand{\arraystretch}{1.25}
\begin{tabular}{r|r r@{\qquad}r|r r}
$\theta$ & & $\lambda$ & $\theta$ & & $\lambda$ \\
\cline{1-1}\cline{3-4}\cline{6-6}
$22.4$ & $+60^{\rm s}20$ & \raisebox{-7pt}{$-0^{\rm s}96$} & $42.4$ & $+59^{\rm s}88$ & \raisebox{-7pt}{$-0^{\rm s}62$} \\
$25.0$ & $59.24$ & \raisebox{-7pt}{$+1.27$} & $48.3$ & $59.24$ & \raisebox{-7pt}{$+3.14$} \\
$28.0$ & $60.51$ & \raisebox{-7pt}{$-0.58$} & $56.2$ & $62.38$ & \raisebox{-7pt}{$-1.06$} \\
$32.9$ & $59.93$ & \raisebox{-7pt}{$-0.85$} & $61.6$ & $61.32$ & \raisebox{-7pt}{$+0.27$} \\
$35.9$ & $59.08$ & \raisebox{-7pt}{$+0.06$} & $62.2$ & $61.59$ & \raisebox{-7pt}{$-2.45$} \\
$37.1$ & $59.14$ & \raisebox{-7pt}{$+1.47$} & $66.8$ & $59.14$ & \raisebox{-7pt}{$+1.05$} \\
$40.4$ & $60.63$ & \raisebox{-7pt}{$-0.75$} & $68.0$ & $60.19$ &  \\
$42.4$ & $59.88$ &  &  &  &  \\
\end{tabular}
\endgroup
\end{center}

Also $S=6.00$, $m=\BareRoot{\dfrac{6.00}{13-3}}$. Hieraus der mittlere zu befürchtende Fehler
bei $x$: $0^{\rm s}75$ bei $y$: $0^{\rm s}40$

Die sämmtlichen von mir berechneten 5 Chronometer geben

\begin{center}
\begin{tabular}{ll|l|l}
 & & E.\,med. & Gewicht \\
\hline
Breguet & $x=1890^{\rm s}36$ & $0^{\rm s}75$ & 1.78 \\
Kessels & $1893.29$ & $0.67$ & 2.23 \\
Barraud & $1892.32$ & $0.49$ & 4.16 \\
Engl. 1 & $1892.39$ & $0.43$ & 5.41 \\
---\quad 4 & $1892.52$ & $0.35$ & 8.16 \\
\hline
Mittel & $x=1892.35$ & & 21.74 \\
\end{tabular}
\end{center}

\begin{center}
\begin{tabular}{ll|l|l}
Breguet & $y=494.12$ & $0.40$ & 6.25 \\
Kessels & $493.89$ & $0.36$ & 7.72 \\
Barraud & $493.67$ & $0.26$ & 14.79 \\
Engl. 1 & $493.98$ & $0.29$ & 11.89 \\
---\quad 4 & $494.16$ & $0.24$ & 17.36 \\
\hline
 & $y=493.96$ & & 58.01 \\
\end{tabular}
\end{center}

Uebrigens ist zwar hier in die letzte Columne unter der Ueberschrift Gewicht $\dfrac{1}{\text{Quadr. (E.\,m.)}}$
gesetzt, also als Einheit die Genauigkeit verstanden, wo der mittlere zu befürchtende Fehler $=1^{\rm s}$
ist, so dass also z.\,B. für Altona der mittlere zu befürchtende Fehler $=\dfrac{1^{\rm s}}{\BareRoot{58.01}}=0^{\rm s}13$
wird; inzwischen wird es rathsamer sein, die Zahlen der letzten Columne blos als Verhältnisszahlen
zu betrachten und die absolute Genauigkeit aus den Unterschieden der aus den einzelnen Chronometern
für $x$ und $y$ gefundenen Werthe von den Endresultaten abzuleiten. Inzwischen wird so die
Genauigkeit des Endresultats
\SourceLeaf{464}noch immer etwas grösser scheinen, als sie wirklich ist, da die Zeitbestimmungen in Greenwich, Helgoland und Altona keine \emph{absolute} Genauigkeit haben, und also offenbar, wenn die Anzahl der Chronometer
auch noch so gross wäre, doch immer die aus jener Quelle entsprungenen Fehler in den Endresultaten nachwirken müssen.

Die Längenbestimmung von Bremen kann auf folgende Art gemacht werden. Setzt man die
Länge $=z$ östlich von Helgoland, so giebt die Vergleichung des \textsc{Breguet}sschen Chronometers

\[
-165^{\rm s}52+z
\]


\begin{center}\small
\begin{tabular}{p{0.42\linewidth} r l r}
Also aus der vorhergenden Vergleichung & $z=225.40$ & Gewicht & $\dfrac{1}{1.4}=0.7$ \\[5pt]
aus der folgenden Vergleichung & $z=224.76$ & & $\dfrac{1}{4.5}=0.2$ \\[8pt]
\cline{2-2}\cline{4-4}
 & $225.24$ & & $0.9$ \\
\end{tabular}
\end{center}

Das Gewicht $0.9$ ist noch mit $\dfrac{10}{6.00}$ zu multipliciren. So geben die 5 Chronometer

\begin{center}
\begin{tabular}{l|l|l}
 & & Gewicht \\
\hline
Breguet & $225.24$ & $1.5$ \\
Kessels & $225.84$ & $1.9$ \\
Barraud & $225.39$ & $3.6$ \\
Engl. 1 & $226.04$ & $2.9$ \\
---\quad 4 & $224.86$ & $4.3$ \\
\hline
 & $225.42$ & $14.2$ \\
\end{tabular}
\end{center}

Allein die Länge von Bremen, die hienach gegen Altona $268^{\rm s}54$ westlich ausfällt, bleibt natürlich immer von der Zeitbestimmung in Bremen abhängig und dieser Unterschied scheint mehrere Secunden zu klein zu sein. Nach meinen Dreiecken ist der Ansgariusthurm $273^{\rm s}51$ in Zeit westlich von
Göttingen, also \textsc{Olbers} Observatorium $271^{\rm s}9$.
